% Source: physical pp. 134--135 (printed pp. 111--112). Coverage:
% Problems 1--6 and References 1--9 plus 7a, ending before Chapter 27.
% Figures/tables/footnotes/dividers: none.

\chapterbackmatter{Problems}

\begin{enumerate}
  \item Using the direct technique of Section 26.1 in the case of $N=2$
  supersymmetry, find the supersymmetry transformation rules of the
  massive field supermultiplet with just one Weyl spinor field and its
  conjugate, and two complex scalars.

  \item Calculate the component fields of a time-reversed superfield
  \[
    T^{-1}S(x,\theta)T
  \]
  in terms of the components of the superfield $S(x,\theta)$. What sort
  of superfield do we get by the time reversal of a left-chiral
  superfield? Of a linear superfield?

  \item Consider the $N=1$ supersymmetric theory of a single left-chiral
  superfield $\Phi$. In superfield notation, list all the terms of
  dimensionality 5 involving $\Phi$ and/or $\Phi^*$ that might be added
  to the Lagrangian density.

  \item Consider the theory of three left-chiral scalar superfields
  $\Phi_1$, $\Phi_2$, and $\Phi_3$, with a conventional kinematic term,
  and with a superpotential
  \[
    f(\Phi_1,\Phi_2,\Phi_3)
    = \Phi_1\Phi_3^2+\Phi_2\bigl(\Phi_3^2+a\bigr),
  \]
  where $a$ is a non-zero real constant. Show that this is a theory with
  spontaneously broken supersymmetry. Find the minimum value of the
  potential. Express the field of the goldstino in terms of the
  fermionic components of $\Phi_1$, $\Phi_2$, and $\Phi_3$.

  \item Find all the components of the current superfield for the action
  (26.6.9), in terms of the components of the left-chiral superfields
  $\Phi_n$ and derivatives of the superpotential $f$ and the Kahler
  potential $K$.

  \item Check that the supersymmetry current given by Eqs. (26.7.18) and
  (26.7.10) is related to the $\omega$-component of the superfield
  (26.7.21) by Eq. (26.7.20).
\end{enumerate}

\chapterbackmatter{References}

\begin{enumerate}
  \item \phantomsection\label{ch26-ref-1}
  A.~Salam and J.~Strathdee, \textit{Nucl. Phys.} B\textbf{76}, 477
  (1974). This article is reprinted in \textit{Supersymmetry},
  S.~Ferrara, ed. (North Holland/World Scientific,
  Amsterdam/Singapore, 1987).

  \item \phantomsection\label{ch26-ref-2}
  J.~Wess and B.~Zumino, \textit{Nucl. Phys.} B\textbf{70}, 13 (1974).
  This article is reprinted in \textit{Supersymmetry},
  Ref.~\hyperref[ch26-ref-1]{1}.

  \item \phantomsection\label{ch26-ref-3}
  L.~O'Raifeartaigh, \textit{Nucl. Phys.} B\textbf{96}, 331 (1975).
  This article is reprinted in \textit{Supersymmetry},
  Ref.~\hyperref[ch26-ref-1]{1}.

  \item \phantomsection\label{ch26-ref-4}
  F.~A.~Berezin, \textit{The Method of Second Quantization}
  (Academic Press, New York, 1966).

  \item \phantomsection\label{ch26-ref-5}
  J.~Iliopoulos and B.~Zumino, \textit{Nucl. Phys.} B\textbf{76}, 310
  (1974); S.~Ferrara and B.~Zumino, \textit{Nucl. Phys.} B\textbf{87},
  207 (1975). These articles are reprinted in \textit{Supersymmetry},
  Ref.~\hyperref[ch26-ref-1]{1}.

  \item \phantomsection\label{ch26-ref-6}
  This section follows the approach of S.~Ferrara and B.~Zumino,
  Ref.~\hyperref[ch26-ref-5]{5}.

  \item \phantomsection\label{ch26-ref-7}
  X.~Gr\`acia and J.~Pons, \textit{J. Phys.} A\textbf{25}, 6357 (1992).
  I am grateful to J.~Gomis for suggesting the use of the equation
  matching the coefficients of $\dot q^n$.

  \item[7a.] \phantomsection\label{ch26-ref-7a}
  M.~Grisaru, in \textit{Recent Developments in Gravitation---Carg\`ese
  1978}, M.~L\'evy and S.~Deser, eds. (Plenum Press, New York, 1979):
  577.

  \item \phantomsection\label{ch26-ref-8}
  B.~Zumino, \textit{Phys. Lett.} \textbf{87B}, 203 (1979). This article
  is reprinted in \textit{Supersymmetry},
  Ref.~\hyperref[ch26-ref-1]{1}.

  \item \phantomsection\label{ch26-ref-9}
  P.~Griffiths and J.~Harris, \textit{Principles of Algebraic Geometry}
  (Wiley, New York, 1978): 109. I thank D.~Freed for telling me of this
  application of the general theorem.
\end{enumerate}
